% !TEX program = pdflatex
\documentclass[11pt]{article}

\usepackage[T1]{fontenc}
\usepackage{lmodern}
\usepackage{microtype}
\usepackage[a4paper,margin=28mm,headheight=14pt]{geometry}
\usepackage{amsmath,amssymb,amsthm,mathtools}
\usepackage{enumitem}
\usepackage{xcolor}
\usepackage{xurl}
\usepackage{hyperref}
\usepackage{fancyhdr}

\definecolor{linkblue}{RGB}{24,78,119}
\definecolor{citegreen}{RGB}{36,104,71}
\hypersetup{
  colorlinks=true,
  linkcolor=linkblue,
  citecolor=citegreen,
  urlcolor=linkblue,
  pdftitle={A Dirichlet Polya Inequality for Three-Dimensional Spherical Shells},
  pdfsubject={Review draft},
  pdfkeywords={Polya conjecture, Dirichlet Laplacian, spherical shell, eigenvalue counting}
}

\pagestyle{fancy}
\fancyhf{}
\fancyhead[L]{\small Dirichlet P\'olya inequality for spherical shells}
\fancyhead[R]{\small Review draft}
\fancyfoot[C]{\thepage}
\fancypagestyle{plain}{%
  \fancyhf{}%
  \fancyhead[L]{\small Dirichlet P\'olya inequality for spherical shells}%
  \fancyhead[R]{\small Review draft}%
  \fancyfoot[C]{\thepage}%
}
\fancypagestyle{titlepage}{%
  \fancyhf{}%
  \fancyfoot[C]{\thepage}%
  \renewcommand{\headrulewidth}{0pt}%
}

\numberwithin{equation}{section}
\allowdisplaybreaks[2]
\setlist{itemsep=0.25em,topsep=0.45em}
\setlength{\emergencystretch}{3em}

\newtheorem{theorem}{Theorem}[section]
\newtheorem{proposition}[theorem]{Proposition}
\newtheorem{lemma}[theorem]{Lemma}

\newcommand{\R}{\mathbb R}
\newcommand{\N}{\mathbb N}
\newcommand{\Nzero}{\mathbb N_0}
\newcommand{\repo}[1]{\nolinkurl{#1}}
\newenvironment{proofrecord}
  {\begin{quote}\small\noindent\textbf{Internal proof record.}\ }
  {\end{quote}}

\title{A Dirichlet P\'olya Inequality for\\Three-Dimensional Spherical Shells\\[0.4em]
\large Dependency-compressed review draft}
\author{}
\date{15 July 2026}

\begin{document}

\maketitle
\thispagestyle{titlepage}

\begin{quote}
\small
\textbf{Scope and status.}
This document reorganizes an internally validated proof into a conventional
review manuscript.  It proves a theorem for the class of genuine
three-dimensional spherical shells.  It does not claim the general P\'olya
conjecture, literature novelty, priority, or publication readiness.  The
final theorem assembly and the non-tiling argument are written out.  The long
compact-middle argument is presented as a sequence of precise lemmas with
exact domains, seams, constants, and pointers to the audited detailed proofs.
A journal submission should move those detailed arguments and the
computer-assisted contracts into conventional appendices rather than cite
repository workflow records.
\end{quote}

\begin{abstract}
For \(0<r<R\), let
\[
  A_{r,R}=\{x\in\R^3:r<|x|<R\}.
\]
We prove, with the strict Dirichlet counting convention, that
\[
  N_D(A_{r,R},\Lambda)
  \le L_3|A_{r,R}|\Lambda^{3/2},
  \qquad
  L_3=\frac1{6\pi^2},
  \qquad \Lambda\ge0.
\]
The proof separates variables, converts the radial spectrum into a strict
multiplicity-weighted phase count, and combines analytic endpoint estimates,
finite angular and radial mode inventories, and a rigorously certified final
compact window.  We also give an elementary proof that no genuine spherical
shell tiles \(\R^3\) by congruent rigid-motion copies with disjoint interiors,
under either exact or almost-everywhere coverage.
\end{abstract}

\section{Main results and conventions}

The eigenvalues of the Dirichlet Laplacian on a bounded domain \(\Omega\) are
listed with multiplicity as
\[
  0<\lambda_1^D(\Omega)\le\lambda_2^D(\Omega)\le\cdots.
\]
Throughout we use the strict counting function
\begin{equation}
  N_D(\Omega,\Lambda)
  :=\#\{j:\lambda_j^D(\Omega)<\Lambda\}.
  \label{eq:strict-count}
\end{equation}
This convention matters at every spectral wall.  We write
\(\N=\{1,2,3,\ldots\}\) and \(\Nzero=\{0,1,2,\ldots\}\).

\begin{theorem}[Spectral inequality]
\label{thm:spectral}
For every \(0<r<R\) and every \(\Lambda\ge0\),
\begin{equation}
  \boxed{
  N_D(A_{r,R},\Lambda)
  \le \frac{2}{9\pi}(R^3-r^3)\Lambda^{3/2}
  =L_3|A_{r,R}|\Lambda^{3/2},
  \qquad L_3=\frac1{6\pi^2}.}
  \label{eq:main-theorem}
\end{equation}
\end{theorem}

\begin{theorem}[Non-tiling]
\label{thm:nontiling}
For every \(0<r<R\), no family of congruent rigid-motion copies of
\(A_{r,R}\) has pairwise-disjoint interiors and covers \(\R^3\) exactly or
up to a three-dimensional Lebesgue-null set.  The same conclusion holds when
coverage is formulated using closed shells.
\end{theorem}

The two theorems concern literally the same complete shell family.
Theorem~\ref{thm:nontiling} is logically independent of the spectral proof.

\section{Separation of variables and strict phase counting}

We first normalize the outer radius to one and write
\[
  A_{\rho,1}=\{x\in\R^3:\rho<|x|<1\},
  \qquad 0<\rho<1.
\]

\begin{proposition}[Exact separated spectrum]
\label{prop:separated-spectrum}
Under spherical-harmonic decomposition and the unitary radial substitution
\(u(r)=rR(r)\), the Dirichlet Laplacian is the orthogonal sum
\begin{equation}
  -\Delta_{A_{\rho,1}}^D
  \simeq
  \bigoplus_{\ell=0}^{\infty}
  \bigoplus_{m=-\ell}^{\ell}L_\ell,
  \qquad
  L_\ell=-\frac{d^2}{dr^2}+\frac{\ell(\ell+1)}{r^2},
  \label{eq:operator-sum}
\end{equation}
where \(L_\ell\) acts on \(H^2(\rho,1)\cap H_0^1(\rho,1)\).  Each radial
block has a positive, simple, complete spectrum.  The full resolvent is
compact because
\begin{equation}
  \|(L_\ell+1)^{-1}\|
  \le \frac1{1+\ell(\ell+1)}\longrightarrow0.
  \label{eq:resolvent-tail}
\end{equation}

Put \(\nu=\ell+\tfrac12\).  The positive eigenfrequencies of the
\(\ell\)-th block are exactly the positive zeros of
\begin{equation}
  F_{\nu,\rho}(k)
  =J_\nu(\rho k)Y_\nu(k)-Y_\nu(\rho k)J_\nu(k).
  \label{eq:bessel-determinant}
\end{equation}
For \(K>0\), choose the continuous Bessel phase by
\[
  J_\nu(x)+iY_\nu(x)=M_\nu(x)e^{i\theta_\nu(x)},
  \qquad \theta_\nu(0+)=-\frac\pi2,
\]
and define
\begin{equation}
  \Psi_{\nu,\rho}(K)=\theta_\nu(K)-\theta_\nu(\rho K),
  \qquad
  \gamma_{\rho,K}(\nu)=\frac{\Psi_{\nu,\rho}(K)}\pi.
  \label{eq:phase-count}
\end{equation}
The function \(K\mapsto\Psi_{\nu,\rho}(K)\) is strictly increasing from
zero to infinity, and
\begin{equation}
  F_{\nu,\rho}(K)=0
  \quad\Longleftrightarrow\quad
  \Psi_{\nu,\rho}(K)\in\pi\N.
  \label{eq:phase-levels}
\end{equation}
Consequently,
\begin{equation}
  N_D(A_{\rho,1},K^2)
  =\sum_{\nu_\ell<K}2\nu_\ell
  [\gamma_{\rho,K}(\nu_\ell)]_<,
  \qquad \nu_\ell=\ell+\frac12,
  \label{eq:strict-phase-sum}
\end{equation}
where
\begin{equation}
  [x]_<:=\#\{n\in\N:n<x\}
  =\max\{0,\lceil x\rceil-1\}.
  \label{eq:strict-bracket}
\end{equation}
In particular, \([m]_< =m-1\) when \(m\in\N\).  If several angular
channels share an eigenfrequency, all equal-threshold modes are excluded at
that wall and enter above it with the sum of their multiplicities.
\end{proposition}

The phase convention and special-function identities used here are
standard; see \cite{DLMF,FLPSphase}.  The exact reconstruction statement is
\repo{state/lemma_packets/SHELL-sturm-liouville-completeness.md}.  The
accepted proof evidence consists of
\repo{rounds/polya-main/round_004/responses/sturm-liouville-incumbent.md},
\repo{rounds/polya-main/round_004/reviews/sturm-liouville-clean-room.md}, and
\repo{rounds/polya-main/round_004/reviews/sturm-liouville-adversarial.md}.

\begin{proposition}[Phase-to-lattice reduction]
\label{prop:phase-to-lattice}
Let \(K>0\).  For \(a>0\), set
\begin{equation}
  G_a(z)=\frac1\pi
  \left(\sqrt{a^2-z^2}-z\arccos\frac za\right),
  \quad 0\le z\le a,
  \label{eq:model-action}
\end{equation}
and \(G_a(z)=0\) for \(z>a\).  Put \(\mu=\rho K\) and
\begin{equation}
  A_{\rho,K}(z)=G_K(z)-G_\mu(z).
  \label{eq:shell-action}
\end{equation}
For \(0\le z<\mu\), define
\[
  H_\mu(z)=
  \frac{3\mu^2+2z^2}{24\pi(\mu^2-z^2)^{3/2}},
\]
and define the auxiliary function---denoted \(F_\mu\) in the annulus
source, but written \(\mathcal F_\mu\) here to distinguish it from the
Bessel determinant---by
\[
  \mathcal F_\mu(z)=
  \begin{cases}
    \max\{G_\mu(z)-H_\mu(z),-1/4\},&0\le z<\mu,\\
    -1/4,&z\ge\mu.
  \end{cases}
\]
Finally, \(\Phi_{K,\mu}(z)=G_K(z)-G_\mu(z)=A_{\rho,K}(z)\) on
\(0\le z\le\mu\).

The global phase enclosure \cite[Theorem~5.1 and Lemma~5.2]{FLPSannuli}
gives, for \(0\le\nu\le K\),
\begin{equation}
  0<\gamma_{\rho,K}(\nu)
  <G_K(\nu)-\mathcal F_\mu(\nu)
  \le G_K(\nu)+\frac14,
  \label{eq:global-phase-bound}
\end{equation}
and, when \(\nu\le\mu\), the sharper two-sided action enclosure
\begin{equation}
  \Phi_{K,\mu}(\nu)
  <\gamma_{\rho,K}(\nu)
  <\Phi_{K,\mu}(\nu)+\frac14.
  \label{eq:action-phase-bound}
\end{equation}
Set
\[
  s_\mu(\nu)=
  \begin{cases}
    \min\{H_\mu(\nu),1/4\},&0\le\nu<\mu,\\
    1/4,&\mu\le\nu\le K,
  \end{cases}
  \qquad
  P_{\rho,K}(\nu)=A_{\rho,K}(\nu)+s_\mu(\nu).
\]
Combining the two phase bounds gives
\[
  \gamma_{\rho,K}(\nu)<P_{\rho,K}(\nu)
  \le A_{\rho,K}(\nu)+\frac14.
\]
Thus the strict spectral count is bounded by the optimized proxy, and in
particular by the coarse proxy through the explicit chain
\[
  N_D(A_{\rho,1},K^2)
  \le
  \sum_{\substack{\ell\ge0\\ \nu_\ell<K}}
  2\nu_\ell[P_{\rho,K}(\nu_\ell)]_<
  \le S_{\rho,K}.
\]
Here
\begin{equation}
  S_{\rho,K}
  =\sum_{\substack{\ell\ge0\\ \nu_\ell<K}}2\nu_\ell
  \left[A_{\rho,K}(\nu_\ell)+\frac14\right]_<.
  \label{eq:coarse-proxy}
\end{equation}
The continuous target is exact \cite{FLPSballs}:
\begin{equation}
  \int_0^\infty2zA_{\rho,K}(z)\,dz
  =\frac{2}{9\pi}(1-\rho^3)K^3.
  \label{eq:action-integral}
\end{equation}
\end{proposition}

The remaining task is to show that strict floor losses, phase slack, and the
half-integer multiplicity mesh do not exceed the target in
\eqref{eq:action-integral}.  The exact reconstruction contracts are
\repo{state/lemma_packets/SHELL-phase-enclosures.md} and
\repo{state/lemma_packets/SHELL-weighted-lattice-fractional.md}.  Accepted
proof evidence is supplied by the Round~2 phase incumbent, clean-room proof,
and adversarial review, and by the corresponding Round~3 low-interface
records:
\begin{proofrecord}
\repo{rounds/polya-main/round_002/responses/phase-enclosure-incumbent.md};
\repo{rounds/polya-main/round_002/reviews/phase-enclosure-clean-room.md};
\repo{rounds/polya-main/round_002/reviews/phase-enclosure-adversarial.md};
\repo{rounds/polya-main/round_003/responses/low-interface-fixed-rho-incumbent.md};
\repo{rounds/polya-main/round_003/reviews/low-interface-fixed-rho-clean-room.md};
\repo{rounds/polya-main/round_003/reviews/low-interface-fixed-rho-adversarial.md}.
\end{proofrecord}

For later use, write
\begin{equation}
  W(\rho,K)=\frac{2}{9\pi}(1-\rho^3)K^3.
  \label{eq:weyl-target}
\end{equation}

\section{The two ratio endpoints}

Define
\begin{equation}
  \omega_0=\frac{\sqrt3}{2\pi}-\frac16,
  \qquad
  C_*=\frac12+\frac{8\sqrt2}{15},
  \qquad
  \rho_*=\frac{\omega_0}{1+2C_*}.
  \label{eq:endpoint-constants}
\end{equation}

\begin{lemma}[Small-hole endpoint]
\label{lem:small-hole}
For \(0<\rho\le\rho_*\) and \(K\ge0\),
\begin{equation}
  N_D(A_{\rho,1},K^2)\le W(\rho,K).
  \label{eq:small-hole-theorem}
\end{equation}
\end{lemma}

\begin{proof}
The multiplicity identity rewrites \eqref{eq:coarse-proxy} as a sum of
shifted tails, one for each \(r\in\Nzero\) and beginning at
\(x_r=r+\tfrac12\).  The ordinary-floor tail majorant is legitimate because
\([x]_<\le\lfloor x\rfloor\) for every \(x\ge0\).  Every tail beginning at
or above the inner interface \(\mu=\rho K\) is controlled by the convex
high-angular theorem of \cite{FLPSballs}.

The separately proved low-interface estimate, using the concave and convex
floor-sum estimates from \cite{FLPSannuli}, controls a tail beginning below
\(\mu\) whenever
\begin{equation}
  0<\rho<\omega_0,
  \qquad K(\omega_0-\rho)\ge C_*.
  \label{eq:small-hole-condition}
\end{equation}
If \(\mu\le\tfrac12\), then every \(x_r\ge\tfrac12\ge\mu\), so only the
high-angular case occurs.  If \(\mu>\tfrac12\) and
\(0<\rho\le\rho_*\), then
\begin{equation}
  K(\omega_0-\rho)
  >\frac{\omega_0-\rho}{2\rho}
  \ge\frac{\omega_0-\rho_*}{2\rho_*}=C_*.
  \label{eq:small-hole-closure}
\end{equation}
Hence every tail is controlled.  Summing the tail integrals gives the target
bound corresponding to \eqref{eq:action-integral}, which proves
\eqref{eq:small-hole-theorem}.  The face \(\rho=\rho_*\) is included.  At
\(K=0\), both sides vanish.
\end{proof}

The exact statement and shifted-tail inequalities are in
\repo{state/lemma_packets/SHELL-rho-zero-endpoint.md}.  The proof and its
independent checks are
\repo{rounds/polya-main/round_007/responses/small-hole-incumbent.md},
\repo{rounds/polya-main/round_007/reviews/small-hole-clean-room.md}, and
\repo{rounds/polya-main/round_007/reviews/small-hole-adversarial.md}.

\begin{lemma}[Thin-shell endpoint]
\label{lem:thin-shell}
For \(7/8\le\rho<1\) and \(K\ge0\),
\begin{equation}
  N_D(A_{\rho,1},K^2)\le W(\rho,K).
  \label{eq:thin-shell-theorem}
\end{equation}
\end{lemma}

\begin{proof}[Proof outline with exact closing margins]
Put \(\varepsilon=1-\rho\in(0,1/8]\) and \(a=\varepsilon K\).  The proof
has two inclusive pieces.

On \(0\le a\le\pi/(4\varepsilon)\), the radial form comparison has the
correct upper-count direction.  Strict radial counting gives the number of
active one-dimensional radial levels as
\[
  M=\max\left\{0,\left\lceil\frac a\pi\right\rceil-1\right\}.
\]
Define
\[
  I(a)=\frac{2a^3}{3\pi},
  \qquad
  S_0(a)=\sum_{n=1}^{M}(a^2-n^2\pi^2),
  \qquad
  D(a)=I(a)-S_0(a).
\]
For \(a>\pi\), write \(a/\pi=M+t\), where \(M\ge1\) and
\(0<t\le1\), using \((M,t)=(m-1,1)\) at \(a=m\pi\).  The exact product
deficit satisfies
\begin{equation}
  \frac{D(a)}{\pi^2}
  =\frac{M^2}{2}+M\left(t^2+\frac16\right)+\frac{2t^3}{3},
  \qquad D(a)>\frac25a^2.
  \label{eq:product-deficit}
\end{equation}
The angular-ceiling and quarter-disk cost is at most
\[
  \left(\frac16+\frac\varepsilon4+\frac{\varepsilon^2}{9}\right)a^2.
\]
At \(\varepsilon=1/8\), the remaining reserve is
\begin{equation}
  \frac25-\left(\frac16+\frac1{32}+\frac1{576}\right)
  =\frac{577}{2880}>0.
  \label{eq:product-reserve}
\end{equation}

On \(a\ge\pi/(4\varepsilon)\), let
\(\mathcal A:[0,a]\to[0,T]\), \(T=a/\pi\), be the exact decreasing shell
action from the phase reduction, let
\(R_{\mathcal A}:[0,T]\to[0,a]\) be its decreasing inverse, and put
\(F_{\mathcal A}=R_{\mathcal A}^2\) and
\(g=-F_{\mathcal A}'\).  Splitting at the actual ungridded curvature
interface, shifted-sawtooth Stieltjes integration gives the radial deficit
\[
  D_{\mathrm{rad}}:=
  \int_0^T F_{\mathcal A}(t)\,dt
  -\sum_{\substack{n\ge1\\ n-1/4<T}}F_{\mathcal A}(n-1/4)
\]
and the bound
\begin{equation}
  D_{\mathrm{rad}}
  \ge\frac{\rho^2a^2}{4}-\frac{\pi\rho a}{4},
  \label{eq:radial-deficit}
\end{equation}
whereas strict half-integer angular counting costs at most
\begin{equation}
  E=\left(\frac a\pi+\frac14\right)
  \left(\varepsilon a+\frac{\varepsilon^2}{4}\right).
  \label{eq:angular-cost}
\end{equation}
The relevant lower screen decreases and the upper screen increases on
\(0<\varepsilon\le1/8\).  At the closed endpoint,
\[
  \frac{D_{\mathrm{rad}}}{a^2}\ge\frac{21}{256},
  \qquad
  \frac{E}{a^2}<\frac{193}{4096},
\]
so
\begin{equation}
  D_{\mathrm{rad}}-E>\frac{143}{4096}a^2>0.
  \label{eq:thin-reserve}
\end{equation}
Both arguments include the common face
\(a=\pi/(4\varepsilon)\), all radial and angular integer walls, and the face
\(\rho=7/8\).  Together they cover every \(a\ge0\) and prove
\eqref{eq:thin-shell-theorem}.
\end{proof}

The frozen exact statement is
\repo{state/lemma_packets/SHELL-rho-one-endpoint-round16.md}; that packet is
a contract, not the proof.  The complete argument is in
\repo{rounds/polya-main/round_016/responses/thin-endpoint-incumbent.md}, with
the independent reconstruction
\repo{rounds/polya-main/round_016/reviews/thin-endpoint-clean-room.md} and
adversarial audit
\repo{rounds/polya-main/round_016/reviews/thin-endpoint-adversarial-referee.md}.

\section{The compact-middle theorem}

This section compresses the longest part of the proof into exact regional
lemmas.  Each lemma has an analytic proof or a rigorously certified finite
part in the listed repository record.

Set
\begin{equation}
\begin{aligned}
  \rho_0&=\frac1{\sqrt{337}},
  &\rho_c&=\frac1{1+2\pi},
  &z_\rho&=\frac\pi{1-\rho},\\
  k_m(\rho)&=\sqrt{z_\rho^2+m(m+1)}\quad(m\in\Nzero),
  &d&=\frac{\sqrt{337}}2.
\end{aligned}
  \label{eq:compact-constants}
\end{equation}
To make the upper-frequency wall explicit, set
\[
  C_0=1+\frac{8\sqrt2}{15},
  \qquad
  a_\rho=\frac{2\pi\rho}{1-\rho},
  \qquad
  \eta_\rho=G_1\!\left(\max\left\{\rho,\frac12\right\}\right),
\]
\[
  K_0(\rho)=
  \left(
  \frac{\sqrt{a_\rho}+\sqrt{a_\rho+4\eta_\rho C_0}}
       {2\eta_\rho}
  \right)^2,
  \qquad
  H_0(\rho)=\frac{C_*}{\omega_0-\rho}
  \quad(\rho<\omega_0),
\]
and, equivalently to taking the minimum of all active upper owners,
\[
  U(\rho)=
  \begin{cases}
    \min\{K_0(\rho),H_0(\rho)\},&0<\rho<\omega_0,\\
    K_0(\rho),&\omega_0\le\rho<5/6,\\
    \min\{K_0(\rho),54/(1-\rho)^2\},&5/6\le\rho<1.
  \end{cases}
\]
Thus \(U\) is precisely the upper-frequency owner obtained from the
fixed-ratio, small-hole, and seam high-energy thresholds.  On the final live
interval \(\rho_c\le\rho<39/50\), it is exactly \(K_0(\rho)\).

\begin{lemma}[Primitive compact envelope]
\label{lem:primitive-envelope}
The zero-count, small-hole, fixed-ratio high-energy, and thin-shell estimates
cover every compact-middle point except the exact nonrectangular set
\begin{equation}
  \mathcal D_{16}
  =\left\{\rho_*<\rho<\frac78:\
  L(\rho)<K<U(\rho)\right\},
  \label{eq:d16}
\end{equation}
where
\[
  L(\rho)=\max\left\{\frac1{2\rho},\frac\pi{1-\rho}\right\}.
\]
In particular, the accepted analytic envelope already gives the global
high-frequency theorem
\begin{equation}
  0<\rho<1,\qquad K\ge295^2
  \quad\Longrightarrow\quad
  N_D(A_{\rho,1},K^2)\le W(\rho,K).
  \label{eq:global-high-frequency}
\end{equation}
\end{lemma}

The precursor estimates are collected in
\repo{state/lemma_packets/SHELL-rho-compact.md}.  The exact formula
\eqref{eq:d16}, including its open faces and the definition of \(U\), is
proved in \repo{state/lemma_packets/SHELL-rho-compact-round17.md}.

\begin{lemma}[Finite angular and radial staircase]
\label{lem:finite-staircase}
The strict inequality \(N_D(A_{\rho,1},K^2)<W(\rho,K)\) holds on each of
the following regions:
\begin{enumerate}
  \item \(\rho_c\le\rho\le7/8\) and
        \(z_\rho\le K\le k_2(\rho)\);
  \item \(\rho_c\le\rho\le7/8\) and
        \(k_2(\rho)<K\le k_5(\rho)\);
  \item \(\rho_0<\rho<\rho_c\) and
        \(1/(2\rho)<K\le d\);
  \item \(\rho_c\le\rho\le7/8\) and
        \(k_5(\rho)<K\le k_6(\rho)\).
\end{enumerate}
They follow from the one-dimensional lower bound
\begin{equation}
  q_{\ell,n}^2\ge n^2z_\rho^2+\ell(\ell+1),
  \label{eq:radial-lower-bound}
\end{equation}
where \(q_{\ell,n}\) is the \(n\)-th positive radial eigenfrequency in
angular channel \(\ell\), together with fixed-channel shell-to-ball min--max
comparison, exact Bessel-zero inequalities, exhaustive multiplicity
inventories, and strict Weyl payments at every fixed or moving entry wall.
The external first-positive-zero inputs are from \cite{Lorch1993}; their
conversion to shell exclusions, all higher-zero bounds, and all mode
inventories are internal.  The corresponding residual is
\begin{equation}
\begin{aligned}
  \mathcal D_{19}={}&
  \left\{\rho_*<\rho\le\rho_0,
  \frac1{2\rho}<K<U(\rho)\right\}\\
  &\cup\left\{\rho_0<\rho<\rho_c,
  d<K<U(\rho)\right\}\\
  &\cup\left\{\rho_c\le\rho<\frac78,
  k_6(\rho)<K<U(\rho)\right\}.
\end{aligned}
  \label{eq:d19}
\end{equation}
At \(\rho=\rho_0\), the candidate lower fiber is empty because
\(1/(2\rho_0)=d\).  All frequency inequalities in \eqref{eq:d19} are
strict.
\end{lemma}

The proof-free packets
\repo{state/lemma_packets/SHELL-first-angular-bands-round17-claim.md},
\repo{state/lemma_packets/SHELL-next-angular-staircase-round18-claim.md}, and
\repo{state/lemma_packets/SHELL-two-sided-staircase-round19-claim.md} fix the
exact statements; they are not proofs.  The principal accepted proof records
are:
\begin{proofrecord}
\textbf{Round 17:}
\repo{rounds/polya-main/round_017/responses/analytic-compression-incumbent.md};
\repo{rounds/polya-main/round_017/reviews/analytic-compression-clean-room.md};
\repo{rounds/polya-main/round_017/reviews/analytic-compression-adversarial-referee.md}.

\textbf{Round 18:}
\repo{rounds/polya-main/round_018/responses/angular-staircase-incumbent.md};
\repo{rounds/polya-main/round_018/reviews/angular-staircase-clean-room.md};
\repo{rounds/polya-main/round_018/reviews/angular-staircase-constants.md};
\repo{rounds/polya-main/round_018/reviews/angular-staircase-cross-comparison.md};
\repo{rounds/polya-main/round_018/reviews/angular-staircase-adversarial-referee.md}.

\textbf{Round 19:}
\repo{rounds/polya-main/round_019/responses/high-ratio-k6-incumbent.md};
\repo{rounds/polya-main/round_019/exploration/lower-ratio-route.md};
\repo{rounds/polya-main/round_019/reviews/two-sided-staircase-clean-room.md};
\repo{rounds/polya-main/round_019/reviews/two-sided-staircase-constants.md};
\repo{rounds/polya-main/round_019/reviews/two-sided-staircase-adversarial-referee.md}.
\end{proofrecord}

\begin{lemma}[Combined lower, high-staircase, and optical closure]
\label{lem:combined-closure}
The inequality is strict on both lower components of \eqref{eq:d19} and on
\begin{equation}
  \rho_c\le\rho\le\frac78,
  \qquad z_\rho\le K\le k_{11}(\rho).
  \label{eq:k11-region}
\end{equation}
It also holds for every frequency on
\begin{equation}
  \frac{39}{50}\le\rho<1,
  \qquad K\ge0,
  \label{eq:optical-region}
\end{equation}
and is strict there when \(K>0\).  Exact subtraction leaves
\begin{equation}
  \boxed{
  \mathcal D_{20}
  =\left\{\rho_c\le\rho<\frac{39}{50},\quad
  k_{11}(\rho)<K<U(\rho)=K_0(\rho)\right\}.}
  \label{eq:d20}
\end{equation}
The face \(\rho=39/50\) belongs to \eqref{eq:optical-region},
\(K=k_{11}(\rho)\) belongs to \eqref{eq:k11-region}, and \(K=U(\rho)\)
belongs to the inherited high-frequency theorem.
\end{lemma}

The exact statement is fixed by the proof-free
\repo{state/lemma_packets/SHELL-combined-closure-round20-claim.md}.  The
Round~20 high- and lower-route notes are exploratory precursors, not promoted
evidence.  The complete accepted evidence, including the conservative
cap-\(74\) cell, is
\begin{proofrecord}
\repo{rounds/polya-main/round_020/reviews/combined-closure-clean-room.md};
\repo{rounds/polya-main/round_020/reviews/combined-closure-clean-room-addendum.md};
\repo{rounds/polya-main/round_020/reviews/combined-closure-cross-comparison.md};
\repo{rounds/polya-main/round_020/reviews/combined-closure-adversarial-referee.md}.
\end{proofrecord}

\begin{lemma}[Exact closure of the final residual]
\label{lem:exact-closure}
Two final estimates hold:
\begin{equation}
  \frac7{51}\le\rho\le\frac{39}{50},
  \quad 12\le K\le200
  \quad\Longrightarrow\quad
  N_D(A_{\rho,1},K^2)<W(\rho,K),
  \label{eq:compact-certificate-theorem}
\end{equation}
and
\begin{equation}
  \rho_c\le\rho\le\frac{39}{50},
  \quad K\ge200
  \quad\Longrightarrow\quad
  N_D(A_{\rho,1},K^2)<W(\rho,K).
  \label{eq:aggregate-tail-theorem}
\end{equation}
The exact guards are
\begin{equation}
  \frac7{51}<\rho_c,
  \qquad
  k_{11}(\rho)>12
  \quad(\rho_c\le\rho<1).
  \label{eq:closure-guards}
\end{equation}
The restriction \(\rho<1\) in the second guard is essential.  Define
\begin{equation}
  \mathcal C_{21}=\mathcal D_{20}\cap\{K\le200\},
  \qquad
  \mathcal T_{21}=\mathcal D_{20}\cap\{K>200\}.
  \label{eq:d20-split}
\end{equation}
Every point of \(\mathcal C_{21}\) lies in the rectangle
\eqref{eq:compact-certificate-theorem}, and every point of
\(\mathcal T_{21}\) lies in \eqref{eq:aggregate-tail-theorem}.  The sets are
disjoint and exhaust \(\mathcal D_{20}\), with \(K=200\) assigned only to
the compact subtraction owner.  Defining
\[
  \mathcal D_{21}
  :=\mathcal D_{20}\setminus
  (\mathcal C_{21}\cup\mathcal T_{21}),
\]
we therefore have
\begin{equation}
  \boxed{\mathcal D_{21}=\varnothing.}
  \label{eq:d21-empty}
\end{equation}
\end{lemma}

The exact scoped statement is
\repo{state/lemma_packets/SHELL-exact-d20-closure-round21-scoped.md}, which
is not by itself a proof.  Accepted evidence consists of
\begin{proofrecord}
\repo{rounds/polya-main/round_021/exploration/exact-d20-closure.md};
\repo{rounds/polya-main/round_021/reviews/exact-d20-closure-clean-room.md};
\repo{rounds/polya-main/round_021/reviews/exact-d20-closure-clean-room-domain-addendum.md};
\repo{rounds/polya-main/round_021/reviews/exact-d20-closure-certificates-source-exec-replacement.md};
\repo{rounds/polya-main/round_021/reviews/exact-d20-closure-adversarial-referee-source-exec.md}.
\end{proofrecord}

\subsection*{Certificate boundary for Lemma~\ref{lem:exact-closure}}

The finite part of \eqref{eq:compact-certificate-theorem} is an
outward-rounded Arb ball-arithmetic partition into \(10{,}580\) exact
rational leaves.  It certifies strict proxy inequalities, interfaces,
integer walls, corners, and exact coverage only on
\[
  [7/51,39/50]\times[12,200].
\]
It does not certify Bessel roots; the spectral transfer is
Proposition~\ref{prop:phase-to-lattice}.

For \eqref{eq:aggregate-tail-theorem}, \(1{,}286\) outward-rounded Arb boxes
certify only the base predicates at \(K=200\).  Let \(\mathcal B\) be the
aggregate reserve, write
\(\mathcal B_K=\partial_K\mathcal B\) and
\(\mathcal B_{KK}=\partial_K^2\mathcal B\), and let \(F\) be its
one-variable curvature lower screen, with
the exact definitions (A13)--(A22) in the listed aggregate-tail contract.  On
\(7/51\le\rho\le39/50\), the boxes prove precisely
\[
  \mathcal B(\rho,200)>0,
  \qquad
  \mathcal B_K(\rho,200)>0,
  \qquad
  F(\rho)>0,
\]
together with the stated domain guards.  The universal conclusion for
\(K>200\) is analytic: on the same ratio interval and for every \(K\ge200\),
\[
  \mathcal B_{KK}(\rho,K)>F(\rho)>0,
\]
and hence
\begin{equation}
\begin{aligned}
  \mathcal B_K(\rho,K)
  &=\mathcal B_K(\rho,200)
    +\int_{200}^{K}\mathcal B_{KK}(\rho,s)\,ds>0,\\
  \mathcal B(\rho,K)
  &=\mathcal B(\rho,200)
    +\int_{200}^{K}\mathcal B_K(\rho,s)\,ds>0.
\end{aligned}
  \label{eq:aggregate-propagation}
\end{equation}
The two integration premises
\(\mathcal B(\rho,200)>0\) and \(\mathcal B_K(\rho,200)>0\) make these
integrations strict, while \(F(\rho)>0\) supplies the curvature lower bound;
positivity of \(\mathcal B_{KK}\) alone would not suffice.  Thus no
finite-box decision is extrapolated beyond its certified domain.  The exact contracts are
\repo{state/certificate_contracts/ROUND21-compact-proxy-contract.md} and
\repo{state/certificate_contracts/ROUND21-aggregate-tail-contract.md}.

\begin{proposition}[Compact-middle theorem]
\label{prop:compact-middle}
For every \(\rho_*<\rho<7/8\) and \(K\ge0\),
\begin{equation}
  N_D(A_{\rho,1},K^2)\le W(\rho,K).
  \label{eq:compact-middle-theorem}
\end{equation}
\end{proposition}

\begin{proof}
The primitive cover of Lemma~\ref{lem:primitive-envelope}, the successive
closed staircase regions of Lemma~\ref{lem:finite-staircase}, and the three
regions of Lemma~\ref{lem:combined-closure} leave exactly
\(\mathcal D_{20}\).  Lemma~\ref{lem:exact-closure} partitions and closes
that entire final residual.  Equation~\eqref{eq:d21-empty} says that no
uncovered point remains.  All shared faces have a unique declared
subtraction owner, while overlapping theorem domains merely give redundant
proofs.  This proves \eqref{eq:compact-middle-theorem}.
\end{proof}

\section{Proof of the spectral theorem}

We now assemble the regional results without importing any new estimate.
First, elementary bounds give
\[
  3\sqrt3>5>\pi,
\]
and hence \(\omega_0>0\).  Also \(\sqrt3<2\) and \(\pi>3\) give
\(\omega_0<1/6\).  Since \(C_*>0\),
\begin{equation}
  0<\rho_*<\omega_0<\frac16<\frac78.
  \label{eq:ratio-ordering}
\end{equation}
Therefore
\begin{equation}
  (0,1)
  =(0,\rho_*]\mathbin{\dot\cup}
  (\rho_*,7/8)\mathbin{\dot\cup}
  [7/8,1).
  \label{eq:ratio-partition}
\end{equation}

Fix \(0<\rho<1\).  At \(K=0\), positivity in
Proposition~\ref{prop:separated-spectrum} gives
\begin{equation}
  N_D(A_{\rho,1},0)=0=W(\rho,0).
  \label{eq:k-zero}
\end{equation}
Let \(K>0\).  If \(0<\rho\le\rho_*\), apply
Lemma~\ref{lem:small-hole}.  If \(\rho_*<\rho<7/8\), apply
Proposition~\ref{prop:compact-middle}.  If \(7/8\le\rho<1\), apply
Lemma~\ref{lem:thin-shell}.  The three cases in
\eqref{eq:ratio-partition} are disjoint and exhaustive, so
\begin{equation}
  N_D(A_{\rho,1},K^2)
  \le\frac{2}{9\pi}(1-\rho^3)K^3
  \qquad(K\ge0).
  \label{eq:unit-shell-theorem}
\end{equation}

Now
\[
  \omega_3=\frac{4\pi}{3},
  \qquad
  L_3=\frac{\omega_3}{(2\pi)^3}=\frac1{6\pi^2},
  \qquad
  |A_{\rho,1}|=\frac{4\pi}{3}(1-\rho^3),
\]
so
\begin{equation}
  L_3|A_{\rho,1}|=\frac{2}{9\pi}(1-\rho^3).
  \label{eq:weyl-normalization}
\end{equation}
Taking \(K=\sqrt\Lambda\) in \eqref{eq:unit-shell-theorem} proves the
unit-shell form of Theorem~\ref{thm:spectral} for every \(\Lambda\ge0\).

For arbitrary \(0<r<R\), put \(\rho=r/R\).  The unitary dilation
\[
  (D_Ru)(x)=R^{-3/2}u(x/R)
\]
maps \(L^2(A_{\rho,1})\) onto \(L^2(A_{r,R})\), and its Dirichlet form
scales by \(R^{-2}\).  Hence
\[
  \lambda_j^D(A_{r,R})=R^{-2}\lambda_j^D(A_{\rho,1}),
\]
and the strict threshold transforms as
\begin{equation}
  N_D(A_{r,R},\Lambda)
  =N_D(A_{\rho,1},R^2\Lambda).
  \label{eq:dilation-count}
\end{equation}
Applying the unit-shell result at \(R^2\Lambda\),
\[
\begin{aligned}
  N_D(A_{r,R},\Lambda)
  &\le\frac{2}{9\pi}(1-\rho^3)(R^2\Lambda)^{3/2}\\
  &=\frac{2}{9\pi}(R^3-r^3)\Lambda^{3/2}\\
  &=L_3|A_{r,R}|\Lambda^{3/2}.
\end{aligned}
\]
This proves Theorem~\ref{thm:spectral} for every \(R>0\), including
\(R<1\).

\section{Proof that spherical shells do not tile}

We prove Theorem~\ref{thm:nontiling} independently of the spectral argument.
Set \(\alpha=r/R\in(0,1)\) and scale by \(1/R\).  It suffices to consider
\[
  T=A_{\alpha,1}=\{x:\alpha<|x|<1\}.
\]
Every rigid motion has the form \(x\mapsto Qx+a\), with \(Q\) orthogonal.
Because \(T\) is radial, \(Q(T)=T\), so every congruent copy is a translate
\[
  T_a=a+T=\{x:\alpha<|x-a|<1\}.
\]

Assume for contradiction that \(\{T_a:a\in\mathfrak C\}\) has
pairwise-disjoint interiors and covers almost everywhere.  Duplicate centers
are impossible because they give the same nonempty open tile.

Fix a unit vector \(e\) and put
\[
  c=\frac{1+\alpha}{2}e,
  \qquad
  \delta=\frac{1-\alpha}{2}.
\]
Write \(B(x,s)=\{y\in\R^3:|y-x|<s\}\) for the open Euclidean ball.
Then \(B(c,\delta)\subset T\).  Therefore \(T_a\) contains
\(B(a+c,\delta)\), and disjointness implies
\begin{equation}
  |a-b|\ge1-\alpha
  \qquad(a\ne b).
  \label{eq:center-separation}
\end{equation}
Thus the centers are locally finite: a bounded set cannot contain infinitely
many points with fixed positive separation.  They are also countable, since
\[
  \mathfrak C
  =\bigcup_{m\ge1}(\mathfrak C\cap\overline B(0,m))
\]
is a countable union of finite sets.

Each tile boundary is a union of two spheres and is three-dimensional null.
The countable union
\[
  Z=\bigcup_{a\in\mathfrak C}\partial T_a
\]
is therefore null.  Exact open-shell coverage is already a special case of
almost-everywhere coverage.  If closed shells cover exactly or almost
everywhere, removing \(Z\) shows that their open interiors still cover almost
everywhere.  Hence it is enough to assume
\begin{equation}
  \left|\R^3\setminus\bigcup_{a\in\mathfrak C}T_a\right|=0.
  \label{eq:ae-coverage}
\end{equation}

Choose one tile \(T_{a_0}\) and its outer sphere
\[
  \Sigma=\{x:|x-a_0|=1\}.
\]
Fix \(p\in\Sigma\) and put \(u=p-a_0\).  For \(n\ge1\), let
\[
  U_n=B\left(p+\frac2n u,\frac1n\right).
\]
Every \(U_n\) lies outside \(T_{a_0}\), is open, and has positive volume.
By \eqref{eq:ae-coverage}, choose \(x_n\in U_n\cap T_{a_n}\) with
\(a_n\ne a_0\).  Since \(x_n\to p\) and \(|x_n-a_n|<1\), the centers
\(a_n\) eventually lie in a fixed bounded set.  Local finiteness lets us
pass to a subsequence with one fixed center \(a\ne a_0\).  Thus
\(p\in\overline{T_a}\).

The point \(p\) is not in \(T_a\).  Otherwise some
\(B(p,\varepsilon)\) lies in \(T_a\).  Choose
\[
  0<t<\min\{\varepsilon,1-\alpha\}.
\]
Then \(p-tu\in T_{a_0}\cap T_a\), contradicting disjoint interiors.
Therefore every \(p\in\Sigma\) belongs to the boundary of a neighboring
tile.

If \(p\in\Sigma\cap\partial T_a\), then
\(|p-a|\in\{\alpha,1\}\), hence \(|a-a_0|\le2\).  Only finitely many
centers satisfy this bound.  Consequently \(\Sigma\) is covered by finitely
many neighboring inner or outer boundary spheres.

Translate so that \(a_0=0\).  The intersection of \(\Sigma\) with a
neighboring sphere \(\{x:|x-a|=s\}\), \(s\in\{\alpha,1\}\), satisfies
\[
  |x|^2=1,
  \qquad |x-a|^2=s^2,
\]
and hence
\begin{equation}
  2a\cdot x=|a|^2+1-s^2.
  \label{eq:sphere-plane}
\end{equation}
Because \(a\ne0\), \eqref{eq:sphere-plane} is a genuine plane.  Its
intersection with \(\Sigma\) is empty, a tangent point, or a circle, and
therefore has zero surface measure on \(\Sigma\).  A finite union of such
sets cannot cover \(\Sigma\), whose area is \(4\pi\).

There is no coincident-sphere exception.  A neighboring outer sphere equals
\(\Sigma\) only when it has the same center, which is the original tile or a
forbidden duplicate.  A neighboring inner sphere cannot equal \(\Sigma\)
because \(\alpha<1\).  We have reached a contradiction.  Scaling back proves
Theorem~\ref{thm:nontiling}.

\appendix

\section{Logical dependency and evidence boundary}

The proof has four logically distinct layers:
\begin{enumerate}
  \item \textbf{Exact analysis:} operator decomposition, phase
        representation, strict counting, weighted-lattice identities,
        endpoint estimates, and analytic propagation.
  \item \textbf{Finite mode analysis:} explicit radial and angular
        inventories and exact Weyl payments through \(k_{11}\).
  \item \textbf{Certified finite computation:} only the rectangle and
        \(K=200\) base predicates stated after
        Lemma~\ref{lem:exact-closure}.
  \item \textbf{Elementary geometry:} the non-tiling proof, which uses none
        of the spectral work.
\end{enumerate}

The legacy numerical Bessel experiments are not proof inputs.  External
Lorch inequalities \cite{Lorch1993} are used only for explicitly stated
first-positive-zero bounds; the shell variational comparison, radial-index
preservation, multiplicity counts, and Weyl payments are internal arguments.
The audited environment reconstructed the exact specializations from the
published statement but did not have access to the article's full proof, so
an external submission must either cite that theorem conventionally or
replace those bounds by self-contained proofs.  The relevant source cards
are \repo{sources/lorch_bessel_zeros.md},
\repo{sources/round19_bessel_zero_bounds.md}, and
\repo{sources/round20_bessel_zero_bounds.md}.

All repository paths in this draft are relative to the repository root.  They
are printed as paths, rather than as external hyperlinks, because the audited
snapshot at local commit \texttt{d89b609} has not yet been published.  A
review copy must therefore be distributed together with a repository checkout
containing that commit; without the checkout, the internal records should be
treated as unavailable.  The final graph and validation records are
\repo{state/proof_obligations.yml},
\repo{state/last_validation_report.md}, and
\repo{rounds/polya-main/round_022/reviews/derived-state-final-audit.md}.

\section{Work required before external submission}

This draft is mathematically complete only relative to the listed detailed
lemmas and certified supplements.  Before external submission it should be
expanded as follows:
\begin{enumerate}
  \item move the full proofs of Propositions~\ref{prop:separated-spectrum}
        and \ref{prop:phase-to-lattice} into analytic appendices;
  \item replace internal packet language in
        Lemmas~\ref{lem:small-hole}--\ref{lem:exact-closure} with conventional
        numbered propositions and complete proofs;
  \item print the Round~17--20 mode inventories, zero exclusions,
        fixed and moving payments, and the corrected cap-\(74\) argument;
  \item provide a computer-assisted-proof supplement containing the exact
        algorithms, rational partitions, environment, hashes, digest, and
        replay instructions;
  \item give conventional bibliographic citations for all external inputs,
        respecting their exact statement-level scopes;
  \item run an independent human line-by-line reconstruction; and
  \item perform a current literature comparison before using any novelty or
        priority language.
\end{enumerate}
These tasks concern exposition, external verification, and scholarly scope.
They do not broaden Theorem~\ref{thm:spectral} to general domains.

\begin{thebibliography}{9}

\bibitem{FLPSballs}
N.~Filonov, M.~Levitin, I.~Polterovich, and D.~A.~Sher,
\newblock P\'olya's conjecture for Euclidean balls,
\newblock \emph{Invent. Math.} \textbf{234} (2023), 129--169.
\newblock \href{https://doi.org/10.1007/s00222-023-01198-1}{doi:10.1007/s00222-023-01198-1}.

\bibitem{FLPSphase}
N.~Filonov, M.~Levitin, I.~Polterovich, and D.~A.~Sher,
\newblock Uniform enclosures for the phase and zeros of Bessel functions and
their derivatives,
\newblock \emph{SIAM J. Math. Anal.} \textbf{56} (2024), no.~6, 7644--7682.
\newblock \href{https://doi.org/10.1137/24M1642032}{doi:10.1137/24M1642032}.

\bibitem{FLPSannuli}
N.~Filonov, M.~Levitin, I.~Polterovich, and D.~A.~Sher,
\newblock P\'olya's conjecture for Dirichlet eigenvalues of annuli,
\newblock \emph{J. Lond. Math. Soc.} \textbf{113} (2026), no.~2, e70425.
\newblock \href{https://doi.org/10.1112/jlms.70425}{doi:10.1112/jlms.70425}.

\bibitem{Lorch1993}
L.~Lorch,
\newblock Some inequalities for the first positive zeros of Bessel
functions,
\newblock \emph{SIAM J. Math. Anal.} \textbf{24} (1993), no.~3, 814--823.
\newblock \href{https://doi.org/10.1137/0524050}{doi:10.1137/0524050}.

\bibitem{DLMF}
NIST Digital Library of Mathematical Functions,
\newblock Chapter 10: Bessel functions,
\newblock \url{https://dlmf.nist.gov/}, accessed 15 July 2026.

\end{thebibliography}

\end{document}
